% SGA 6, Exposición XIII: fragmento español alineado con la fuente.
% Alcance exacto: fr.tex líneas 14951--15065. El comienzo continúa el
% entorno tikzcd y la enumeración abiertos en el fragmento precedente.
\Phi:\underline{\Pic}_{Y/S}\arrow[r]\arrow[u,"\Phi_Y"]&
\underline{\Pic}_{X/S}\arrow[u,"\Phi_X"]
\end{tikzcd}
\]
\emph{con $\Phi'$ y $\Phi_Y$ de tipo finito, se tiene que $\Phi$ es de tipo
finito.}
\end{enumerate}

En efecto, supongamos que $\Phi_T$ es de tipo finito. Sea $\mathscr L$ una
familia sobre las fibras de $X/S$, acotada por un $L$ sobre un $X_R$. La
familia $\mathscr L_T=\mathscr L\times_S T$ sobre las fibras de $X_T/T$,
definida de manera evidente, está acotada por la imagen inversa de $L$ sobre
$X_{R\times T}$. Si $\Phi_T^{-1}(\mathscr L_T)$ está acotada por un $M$
sobre un $Y_Q$, entonces $M$ acota también $\Phi^{-1}(\mathscr L)$. El
recíproco de (i) y la afirmación (ii) se demuestran de manera análoga.

\textbf{Teorema 3.5.} \emph{Sean $S$ un esquema noetheriano, $X$ e $Y$ dos
$S$-esquemas propios y $f:X\to Y$ un morfismo. Si $f$ es sobreyectivo,
entonces el morfismo inducido entre los funtores de Picard}
\[
f^*:\underline{\Pic}_{Y/S}\longrightarrow\underline{\Pic}_{X/S}
\]
\emph{es de tipo finito.}

En efecto, según 3.4 (i), se puede suponer evidentemente que $S$ es reducido
e irreducible. Existe entonces un abierto no vacío $U$ de $S$ tal que
$\underline{\Pic}_{Y/S}$ y $\underline{\Pic}_{X/S}$ son representables y la
restricción de $f^*$ es cuasi-afín (XII 1.1), y por tanto de tipo finito por
3.2 (i). Dotando a $T=S-U$ de la estructura reducida inducida, por
recurrencia noetheriana se puede suponer que el morfismo
\[
\underline{\Pic}_{Y_T/T}\longrightarrow\underline{\Pic}_{X_T/T}
\]
inducido por $f^*$ es de tipo finito; de aquí resulta la afirmación, gracias
a 3.4 (i) aplicado a $U\amalg T\to S$.

\textbf{Teorema 3.6.} \emph{Sean $S$ un esquema noetheriano, $X$ un
$S$-esquema propio y $\underline{\Pic}_{X/S}$ el funtor de Picard. Entonces,
para todo entero $n\ne0$, el homomorfismo de potencia $n$-ésima}
\[
\Phi_n:\underline{\Pic}_{X/S}\longrightarrow\underline{\Pic}_{X/S}
\]
\emph{es un morfismo de tipo finito.}

En efecto, sea $f:X'\to X$ una presentación de Chow: $X'$ es proyectivo
sobre $S$ y $f$ es sobreyectivo. Si se admite el teorema para $X'$, el
teorema para $X$ resulta de 3.5 y 3.4 (ii). Se puede suponer, por tanto, que
$X/S$ es proyectivo.

Existen un abierto no vacío $U$ de $S$, un esquema íntegro $S'$ y un
morfismo finito y sobreyectivo $S'\to U$ tales que toda componente
irreducible $X'_i$ de $X_{S'}$, dotada de la estructura reducida inducida,
es íntegra sobre $S'$ (XII 1.3). Dotando a $T=S-U$ de la estructura reducida
inducida, por recurrencia noetheriana se puede admitir el teorema para
$X_T/T$. Si se admite para todos los $X_i/S'$, resulta para
$\coprod X_i/S'$ por el lema elemental siguiente, y luego para
$X_{S'}/S'$ en virtud de 3.5 y 3.4 (ii); finalmente resulta para $X/S$
gracias a 3.4 (i).

\textbf{Lema 3.7.} \emph{Sean $X=\coprod X_i$ de tipo finito sobre un
esquema noetheriano $S$ y $\mathcal F$ una familia de clases de haces
coherentes sobre las fibras de $X/S$. Para que $\mathcal F$ sea acotada, es
necesario y suficiente que, para todo $i$, la familia $\mathcal F_i$
inducida sobre las fibras de $X_i/S$ sea acotada.}

En la demostración de 3.6, la cuestión se ha reducido al caso en que $X$ es
proyectivo e íntegro sobre $S$; dotemos a $X/S$ de un haz muy amplio
$\mathcal O_X(1)$. Sean $\mathcal L$ una familia acotada de clases de haces
invertibles sobre las fibras de $X/S$ y $L$ un haz invertible sobre un
$X_T$, con $T$ de tipo finito sobre $S$, que acota $\mathcal L$. Después de
descomponer $T$ en estratos, se puede suponer que $L$ es plano sobre $T$ y,
tras reemplazar $T$ por cada una de sus componentes conexas, que $T$ es
conexo. Entonces el polinomio en $q$ y $m$,
$P(q,m)=\chi(L_t^{\otimes q}(m))$, no depende de $t\in T$ (EGA III 7.9.4).

Sea $\mathcal L_n=\Phi_n^{-1}(\mathcal L)$: $\mathcal L_n$ es la familia de
las clases de haces invertibles $M_K$ tales que la clase de
$M_K^{\otimes n}$ pertenece a $\mathcal L$. Sea
$\chi(M_K^{\otimes p}(m))=\sum_{i=0}^{r}a_i(p)\binom{m+i}{i}$, donde
$a_i(p)$ es un polinomio en $p$ de grado $\leq r-i$. Para $q,m$ variables,
se tiene $\chi(M_K^{\otimes qn}(m))=P(q,m)$, y cada $a_i(qn)$ no depende,
por tanto, del $M_K$ elegido; en consecuencia, cada polinomio $a_i(p)$
tampoco depende de $M_K$. Así pues, la familia $\mathcal L_n$ tiene un único
polinomio de Hilbert $\chi(M_K(m))$ y es, por tanto, acotada en virtud de
2.11.

\textbf{Teorema 3.8.} \emph{Sean $S$ un esquema noetheriano, $Y$ un
$S$-esquema proyectivo provisto de un haz invertible relativamente amplio
$\mathcal O_Y(1)$, $X$ el esquema de ceros de una sección $\varphi$ de
$\mathcal O_Y(1)$ y $f:X\to Y$ el morfismo de inclusión. Se supone que, para
todo punto $s\in S$ y toda componente irreducible $Y_s^i$ de $Y_s$, las
componentes irreducibles de
$X_s\cap Y_s^i=V(\varphi|Y_s^i)$ son de dimensión $\geq2$ (lo que ocurre si
las componentes irreducibles de las fibras de $Y$ son de dimensión
$\geq3$). Entonces el morfismo inducido entre los funtores de Picard}
\[
f^*:\underline{\Pic}_{Y/S}\longrightarrow\underline{\Pic}_{X/S}
\]
\emph{es de tipo finito. Además, si $S$ es íntegro, existe un abierto no
vacío $U$ de $S$ tal que los funtores de Picard restringidos a $U$ son
representables y el morfismo inducido}
\[
\underline{\Pic}_{Y/S}|U\longrightarrow\underline{\Pic}_{X/S}|U
\]
\emph{es cuasi-afín (de tipo finito).}

Demostremos primero la primera afirmación. Como en la demostración de 3.6,
se ve por recurrencia noetheriana que se puede reemplazar $S$ por cualquier
abierto no vacío, y la cuestión se reduce fácilmente al caso en que $S$ es
íntegro e $Y$ es íntegro sobre $S$, reemplazando siempre $X$ por
$X\times_YY'$ cuando se reemplaza $Y$ por $Y'$. Hecho esto, se reduce de la
misma manera al caso en que $Y$ es normal sobre $S$, utilizando, en lugar de
XII 1.3, el lema siguiente.

\textbf{Lema 3.9.} \emph{Sean $S$ un esquema noetheriano íntegro y $X$ un
$S$-esquema íntegro sobre $S$. Entonces existen un abierto no vacío $U$ de
$S$, un esquema íntegro $S'$ finito y fielmente plano sobre $U$, un
$S'$-esquema $X'$ normal e íntegro sobre $S'$, y un $S'$-morfismo finito y
sobreyectivo $X'\to X_{S'}$.}

En efecto, sean $\eta$ el punto genérico de $S$ y $K$ la clausura algebraica
de $k(\eta)$. Es claro que $\Spec(K)$ es el límite del sistema proyectivo
filtrante $(S_\alpha)$ de los esquemas afines íntegros provistos de un
morfismo finito y fielmente plano $S_\alpha\to U_\alpha$, donde $U_\alpha$
es un abierto no vacío de $S$. El normalizado de $X_K$ proviene, por tanto,
de un $S_\alpha$-esquema $X'$ para algún $\alpha$, y el morfismo finito y
sobreyectivo $X'_K\to X_K$ proviene de un $S'_\alpha$-morfismo finito y
sobreyectivo $X'\to X_{S_\alpha}$ (EGA IV 8.8.2, 8.10.5). Finalmente, como
$X'_K$ es normal e íntegro, existe un abierto $S'$ de $S_\alpha$ tal que
$X'_{S'}$ es normal e íntegro sobre $S'$ (EGA IV 9.9.5).

\textbf{Observación 3.10.} Conservando las notaciones y condiciones de 3.9,
supongamos que existe un morfismo resolvente $X'_K\to X_K$ (EGA IV 7.9.1),
por ejemplo si $\dim(X_K)\leq2$, según Abhyankhar. Un razonamiento semejante
muestra entonces que existen un morfismo finito y fielmente plano
$S'\to U$, un esquema $X'$ liso sobre $S'$ y un morfismo propio,
sobreyectivo y birracional $X'\to X_{S'}$.

En la demostración de la primera parte de 3.8, la cuestión se ha reducido al
caso en que $S$ es íntegro e $Y$ es normal e íntegro sobre $S$, y por tanto
$Y$ es íntegro. Si $\varphi=0$, entonces $X=Y$ y la afirmación es trivial.
En caso contrario, tras restringir $S$, se puede suponer que para todo
$s\in S$ se tiene $\varphi_s\ne0$; entonces $\varphi_s$ es
$\mathcal O_{Y_s}$-regular y la dimensión $r$ de $Y_s$ es $\geq3$.

Sea $\mathcal L$ una familia acotada de clases de haces invertibles sobre
las fibras de $X/S$. Sea $M_K$ un haz invertible sobre una fibra geométrica
$Y_K$ de $Y/S$, con polinomio de Hilbert
$\chi(M_K(n))=\sum_{i=0}^{r}a_i\binom{n+i}{i}$. Entonces el polinomio de
Hilbert de $f^*M_K$ es
$\sum_{i=0}^{r-1}a_{i+1}\binom{n+i}{i}$. Cuando la clase de $M_K$ recorre
$(f^*)^{-1}(\mathcal L)$, es decir, cuando la clase de $f^*M_K$ pertenece a
$\mathcal L$, los tres coeficientes $a_r,a_{r-1},a_{r-2}$ permanecen
acotados gracias a 1.13.

Sea $H=\mathcal O_Y(m)$ un múltiplo de $\mathcal O_Y(1)$ que sea muy amplio.
Entonces, en los polinomios de Hilbert
$\chi(M_K(mn))=\sum_{i=0}^{r}a_i^{(m)}\binom{n+i}{i}$ relativos a $H$ de los
$M_K$ en $(f^*)^{-1}(\mathcal L)$, los tres coeficientes
$a_r^{(m)},a_{r-1}^{(m)},a_{r-2}^{(m)}$ permanecen acotados gracias a 2.10.
Por consiguiente, $\mathcal L$ es acotada en virtud de 2.11; de aquí la
primera parte de 3.8.

En cuanto a la segunda parte de 3.8, la representabilidad «genérica» de los
funtores de Picard se estableció en XII 1.2, así como su separación. Por un
razonamiento semejante al de XII 2, utilizando 3.9 y XII 1.1, se reduce la
cuestión al caso en que $Y$ es normal sobre $S$. Ahora bien, para que el
morfismo
$\underline{\Pic}_{Y/S}\to\underline{\Pic}_{X/S}$ sea cuasi-afín, basta
que sea separado y cuasifinito (EGA IV 8.12.8), lo cual se verifica sobre
las fibras geométricas y resulta, por tanto, del lema siguiente.

\textbf{Lema 3.11.} \emph{Sean $k$ un cuerpo algebraicamente cerrado, $Y$ un
$k$-esquema normal y proyectivo provisto de un haz amplio $\mathcal O_Y(1)$,
y $X$ el esquema de ceros de una sección $\varphi$ de $\mathcal O_Y(1)$. Se
supone que $Y$ es de dimensión $\geq3$ (resp. $\geq2$). Entonces el
morfismo}
\[
\Pic_{Y/k}\longrightarrow\Pic_{X/k}
\]
\emph{(resp. el morfismo $\Pic^\tau_{Y/k}\to\Pic^\tau_{X/k}$ (véase 4.1)) es de
tipo finito y tiene un núcleo $N$ que es un grupo finito (no reducido en
general).}

En efecto, se puede suponer evidentemente que $\varphi\ne0$; entonces
$\varphi$ es regular, pues $Y$ es íntegro. El morfismo considerado es de
tipo finito gracias a la primera parte de 3.8 (resp. gracias a 4.7 (iii)).
Por consiguiente, el núcleo $N$ es un grupo algebraico, contenido en
$\Pic^T_{Y/k}$, y por ello es el mismo en los dos casos considerados. Como
$\Pic^T_{Y/k}$ es propio, pues $Y$ es normal (FGA 236-2.13), y como $N$ es
un subgrupo necesariamente cerrado (SGA 3 VI$_B$ 1.9.2), para que $N$ sea
finito basta que no contenga un subgrupo de tipo
$\mu_\ell(=\mathbb Z/\ell\mathbb Z)$, donde
$(\ell,\operatorname{car}(k))=1$.

Consideremos la teoría de Kummer, es decir, el diagrama conmutativo siguiente,
procedente de la sucesión exacta
\[
0\longrightarrow\mu_\ell\longrightarrow\mathbb G_m
\xrightarrow{x\mapsto x^\ell}\mathbb G_m\longrightarrow0
\]
de haces étales sobre $Y$ (resp. sobre $X$):
\[
\begin{tikzcd}[column sep=1.6em,row sep=1.2em]
& & H^1(Y,\mu_\ell)\arrow[r]\arrow[d]&
H^1(Y,\mathbb G_m)\arrow[r,"\ell"]\arrow[d]&
H^1(Y,\mathbb G_m)\arrow[d]\\
H^0(X,\mathbb G_m)\arrow[r,"\ell"]&H^0(X,\mathbb G_m)\arrow[r]&
H^1(X,\mu_\ell)\arrow[r]&H^1(X,\mathbb G_m)\arrow[r,"\ell"]&
H^1(X,\mathbb G_m).
\end{tikzcd}
\]
Por el «teorema 90» de Hilbert, se tiene
$H^1(Y,\mathbb G_m)=\Pic_{Y/k}(k)$ (resp.
$H^1(X,\mathbb G_m)=\Pic_{X/k}(k)$). Como $X$ es propio e íntegro,
$H^0(X,\mathbb G_m)=k^*$; por tanto,
$H^1(X,\mu_\ell)\to H^1(X,\mathbb G_m)$ es inyectivo.

Basta, pues, ver que
$H^1(Y,\mathbb Z/\ell\mathbb Z)\longrightarrow
H^1(X,\mathbb Z/\ell\mathbb Z)$ es inyectivo. Esto resulta de que
$\pi_1(X,\xi)\to\pi_1(Y,\xi)$ es sobreyectivo o, equivalentemente (SGA 1 V
6.9), de que todo recubrimiento étale conexo $\widetilde Y$ de $Y$, de grupo
$\mathbb Z/\ell\mathbb Z$, tiene una restricción
$\widetilde X=X\times_Y\widetilde Y$ que también es conexa. Esta es una
consecuencia bien conocida del teorema de Bertini y del teorema de conexión
de Zariski (SGA 1 X 2.10).

\textbf{Observación 3.12.} Conservemos las condiciones de 3.11. Entonces:
\begin{enumerate}[label=(\roman*)]
\item El núcleo $N$ es un grupo finito \emph{unipotente} (SGA 3 XVII 1.1),
pues no contiene ningún subgrupo de tipo multiplicativo $\mu_n$ para ningún
$n$ (SGA 3 XVII 4.6.1 vi). Esto se demuestra de la misma manera, utilizando,
en lugar de la teoría de Kummer, el hecho siguiente, debido a Grothendieck
(SGA 1 XI, última página): si $X$ es un esquema propio sobre un cuerpo $k$,
con $H^0(X,\mathcal O_X)=k$, existe un isomorfismo de $k$-esquemas en grupos
$\Hom_{\mathrm{gr}}(\mu_n,\underline{\Pic}_{X/k})\simeq
H^1(X,\mathbb Z/n\mathbb Z)$; por consiguiente, si $k$ es algebraicamente
cerrado, se tiene
$\Hom_{\mathrm{gr}}(\mu_n,\underline{\Pic}_{X/k})\simeq
H^1(X,\mathbb Z/n\mathbb Z)$.
\item Si $k$ es de característica cero, entonces el núcleo $N$ es nulo,
pues todo grupo finito unipotente es nulo en ese caso (SGA 3 XVII 1.7). Se
recupera así, aproximadamente, el «Vanishing theorem» de Kodaira en la forma
dada por Mumford [2].
\end{enumerate}

\textbf{Teorema 3.13.} \emph{Sean $S$ un esquema noetheriano y $X$ un
$S$-esquema proyectivo, provisto de un haz invertible relativamente amplio
$H=\mathcal O_X(1)$. Se supone que las fibras geométricas de $X/S$ son
íntegras y de dimensión $r$. Entonces, para una familia $\mathcal L$ de
clases de haces invertibles sobre las fibras de $X/S$, las condiciones
siguientes son equivalentes:}
\begin{enumerate}[label=(\roman*)]
\item \emph{$\mathcal L$ es una familia acotada.}
\item \emph{En los polinomios de Hilbert
$\chi(L_K(n))=\sum_{i=0}^{r}a_i\binom{n+i}{i}$ de los
$L_K\in\mathcal L$, el coeficiente $a_{r-1}$ permanece acotado y el
coeficiente $a_{r-2}$ permanece acotado inferiormente.}
\item \emph{Cuando $L_K$ recorre $\mathcal L$, el grado
$d(L_K)=\langle c_1(L_K)\cdot c_1(H_K)^{r-1}\rangle$ permanece acotado y el
número de intersección
$\langle c_1(L_K)^2\cdot c_1(H_K)^{r-2}\rangle$ permanece acotado
inferiormente.}
\item \emph{Existe un entero $N$ tal que}
\[
H^{-N}<\mathcal L<H^N
\]
\emph{en el sentido de la relación de orden definida por la relación de
amplitud; es decir, para todo $L_K\in\mathcal L$,
$H_K^{\otimes N}\otimes L_K$ y $H_K^{\otimes N}\otimes L_K^{-1}$ son
amplios.}
\end{enumerate}

Demostremos que (i) implica (ii), (iii) y (iv). Supongamos que $\mathcal L$
está acotada por un haz invertible sobre un $X_T$. Entonces $\mathcal L$
satisface (iv) por EGA II 4.6.12. Supongamos que $L$ es plano sobre $T$. El
polinomio en $m$ y $n$, $\chi(L_t^{\otimes m}(n))$, solo depende entonces de
la componente conexa de $T$ (EGA III 7.9.4). Por tanto, $\mathcal L$
satisface (ii) y, teniendo en cuenta 2.1, también (iii).

Observemos también que la familia de las clases de los $H_K$ (resp. de los
$\mathcal O_{X_K}$) es acotada y, por tanto, satisface (iii) (resp. (ii)).

Supongamos que $\mathcal L$ satisface (iv). Entonces
$\langle[Nc_1(H_K)+c_1(L_K)]\cdot c_1(H_K)^{r-1}\rangle>0$, de donde
$d(L_K)>-Nd(H_K)$ y, del mismo modo, $Nd(H_K)>d(L_K)$. Además,
$\langle[Nc_1(H_K)+c_1(L_K)]^2\cdot c_1(H_K)^{r-2}\rangle>0$, de donde
$\langle c_1(L_K)^2\cdot c_1(H_K)^{r-2}\rangle>
-2Nd(L_K)-N^2d(H_K)$. Por tanto, $\mathcal L$ satisface (iii).

Demostremos finalmente que (ii) (resp. (iii)) implica (i). Gracias a 2.10 y
2.3, se puede reemplazar $H$ por un múltiplo; supongamos, pues, que $H$ es
muy amplio. Evidentemente se puede descomponer $S$ en estratos; supongamos
que existe una sección lineal $Y$ de $X$, íntegra sobre $S$, cuyas fibras
son de dimensión $2$. En virtud de 3.8, se puede reemplazar
$X,\mathcal L$ por $Y,\mathcal L|Y$ (teniendo en cuenta 1.7, resp. la fórmula
de proyección, que implica que la condición (ii), resp. (iii), sobre
$X,\mathcal L$ equivale a (ii), resp. (iii), sobre $Y,\mathcal L|Y$);
supongamos, pues, que $r=2$. Si
$\delta=\sup\{d(L_K)\}$ es $>0$, gracias a 2.3 y 2.10 se puede reemplazar
$\mathcal L$ por la familia de las clases de los $L_K(-\delta)$;
supongamos, pues, que $\delta\leq0$.

Si, en los polinomios de Hilbert
$\sum_{i=0}^{2}a_i\binom{n+i}{i}$, $a_1$ permanece acotado y $a_0$
permanece acotado inferiormente, los tres coeficientes permanecen acotados
gracias a 2.3 y 2.7; por tanto, $\mathcal L$ es acotada en virtud de 2.11.
Así, (ii) implica (i). Por otra parte, (iii) implica (ii) gracias a 2.3 y
2.8.
